\section{Introduction}

The assembly of large space structures in orbit has been demonstrated
operationally through programs such as the International Space
Station (ISS), which required over 40 assembly flights across more
than a decade~\cite{nasa_iss_assembly}. However, the scale of
structures envisioned for future missions---kilometer-class solar
power satellites, large interferometric telescope arrays, and
interplanetary propulsion stages---far exceeds what was required for
the ISS and demands a systematic optimization approach to campaign
planning~\cite{wertz_larson}.

The Earth-Sun L4 Lagrange point offers several advantages as an
assembly site for large deep-space vehicles. Its gravitational
stability (the point co-rotates with Earth at 1\,AU from the Sun)
eliminates the stationkeeping penalty associated with more
energetically expensive halo orbits~\cite{belbruno_wbs}.
Continuous solar illumination supports high solar array output,
and the separation from low-Earth orbit (LEO) debris belts reduces
shielding requirements for long-duration operations. The transit
distance of approximately 1\,AU, however, imposes significant
transfer $\Delta v$ requirements (3.8--4.1\,km/s from LEO) and
transit times of 60--120 days, making resupply logistics a primary
cost driver.

Prior work on space logistics optimization has applied
time-expanded network models and dynamic programming to station
resupply~\cite{ho_logistics} and propellant depot
operations~\cite{ho_depot}. This paper extends those foundations
to the on-orbit assembly problem, incorporating a multi-objective
formulation, crew rotation constraints, proximity operations
penalties, and a multi-nation vehicle catalog reflecting the
emerging Artemis Accords partnership structure.

The specific contributions of this work are:
\begin{enumerate}
  \item A dynamic programming formulation over a time-expanded
        network for assembly campaign optimization, with forward
        beam search to manage state-space complexity;
  \item A multi-objective cost function with user-adjustable weights
        over launches, time, and program cost, enabling Pareto
        frontier analysis;
  \item An interactive web-based simulation tool supporting
        parametric spacecraft configuration, real-time playback,
        and one-click PDF report generation.
\end{enumerate}
